\begin{scanpage}{0184}{166}
\[
\left.\begin{aligned}j_x^0&=-\frac c{2\pi}E_{0x}(n_y\sin\gamma+n_z\cos\gamma)e^{i\psi},\\
j_y^0&=\frac c{2\pi}E_{0x}n_x\sin\gamma\,e^{i\psi},\\
j_z^0&=\frac c{2\pi}E_{0x}n_x\cos\gamma\,e^{i\psi},
\end{aligned}\right\}\tag{26.07}
\]
and with H-polarization ($\mathbf H_0\perp yoz$)
\[
\left.\begin{aligned}j_x^0&=0,\\j_y^0&=\frac c{2\pi}H_{0x}n_ze^{i\psi},\\j_z^0&=-\frac c{2\pi}H_{0x}n_ye^{i\psi}.
\end{aligned}\right\}\tag{26.08}
\]
Here $E_{0x}$ and $H_{0x}$ are the electric and magnetic field amplitudes of the incident wave with E-polarization and H-polarization, respectively; $\psi=k(y'\sin\gamma+z'\cos\gamma)$ is the incident wave phase at the point $(x',y',z')$ on the body's surface; $n_x$, $n_y$, $n_z$, are the components of the normal to the surface at the same point.

Furthermore, calculating the vector potential in the far zone on the basis of this current and substituting its values into the equations
\[
\left.\begin{aligned}E_\varphi=-H_\vartheta&=ikA_\varphi,\\E_\vartheta=H_\varphi&=ikA_\vartheta,
\end{aligned}\right\}\tag{26.09}
\]
we find the fringing field. With E-polarization, it equals
\[
\begin{aligned}E_\varphi=-H_\vartheta=\frac{ik}{2\pi}E_{0x}\frac{e^{ikR}}R\int[&n_x\sin\gamma\cos\varphi\\&+(n_y\sin\gamma+n_z\cos\gamma)\sin\varphi]e^{i\Phi}\,dS,
\end{aligned}\tag{26.10}
\]
\[
\begin{aligned}E_\vartheta=H_\varphi=\frac{ik}{2\pi}E_{0x}\frac{e^{ikR}}R\int[&n_x(\sin\gamma\cos\vartheta\sin\varphi-\cos\gamma\sin\vartheta)\\&-(n_y\sin\gamma+n_z\cos\gamma)\cos\varphi\cos\vartheta]e^{i\Phi}\,dS,
\end{aligned}\tag{26.10}
\]
and with H-polarization
\end{scanpage}
